\chapter{The Critique of Indology}

\section{Ghurye Ignored Material Conditions}

Ghurye's Indology is all about culture --- how Hindu culture formed and diffused, and how this produced cultural unity. The first critique is that it ignored \textbf{material conditions}: livelihood, control over resources, class struggle, adaptation. Ghurye gave primacy to \textbf{ideas over matter}.

\begin{ComparisonBox}
\begin{itemize}
\item \textbf{Ghurye} looked at the same historical facts (or mythology) as resulting in \textbf{cultural formation}.
\item \textbf{D. D. Kosambi} (a Marxist) looked at the Aryan invasion as \textbf{class struggle} --- why did the indigenous people become slaves? This focuses on the economic base, on the formation of slavery, not culture.
\end{itemize}
\end{ComparisonBox}

\begin{QuickRevision}
\begin{itemize}
\item Ghurye ignored material conditions; gave primacy to ideas over matter.
\item D. D. Kosambi read the Aryan invasion as class struggle (economic base).
\end{itemize}
\end{QuickRevision}

\section{Book View vs Field View}

Indology provides a \textbf{book view} of Indian society, divorced from the \textbf{field view}. Ghurye looked at caste through the Varna model because that is how the Vedas describe it. But in the field:
\begin{itemize}
\item You do not find Brahmin, Kshatriya, Vaishya, Shudra as four clear groups.
\item You do not find the same fixed hierarchy.
\item Brahmin is not homogeneous --- there is hierarchy within Brahmins.
\item Pollution is not limited to untouchables --- it is found within Brahmins too.
\end{itemize}

\begin{CriticalPoint}
What India is, and what India is projected to be, are two different Indias. The book-view critique is associated with \textbf{M. N. Srinivas}.
\end{CriticalPoint}

\begin{QuickRevision}
\begin{itemize}
\item Indology = book view, divorced from field view.
\item Field reality: no four clear Varnas, no fixed hierarchy, Brahmins heterogeneous, pollution within Brahmins.
\end{itemize}
\end{QuickRevision}

\section{Undue Importance to Brahmins}

Ghurye gave \textbf{undue importance to Brahmins} in the cultural history of India. No one group can be treated as the reason for the formation of Indian culture; and this great status cannot be attributed to Brahmins purely on the basis of \textbf{mythology} (the Aryan invasion theory has no written records).

\begin{DataFact}
In ancient texts you do not find Brahmin dynasties --- dynasties have always been \textbf{Kshatriya (kings)}. There were even times when Brahmins were not considered in dynastic formation (e.g.\ the Mughal empire).
\end{DataFact}

\begin{CriticalPoint}
Ghurye treated the Aryans as Hindu upper-caste Brahmins, and Vedic Aryan Hinduism as the product of Brahmins --- a claim not based on any fact.
\end{CriticalPoint}

\begin{QuickRevision}
\begin{itemize}
\item Ghurye gave undue importance to Brahmins, on the basis of mythology.
\item No Brahmin dynasties; dynasties were Kshatriya (e.g.\ Mughal empire).
\end{itemize}
\end{QuickRevision}

\section{Vedic Hinduism Equated with Indian Culture}

Ghurye \textbf{equated Vedic Hinduism with Indian culture}. In reality Indian culture is \textbf{diverse and composite}, consisting of people professing multiple faiths and ideologies. Islam and Buddhism claim distinct identity --- Buddhism came as a challenge to Brahmanism, so equating Buddhism with Hinduism is wrong.

\begin{QuickRevision}
\begin{itemize}
\item Ghurye equated Vedic Hinduism with Indian culture; but Indian culture is diverse and composite.
\item Islam and Buddhism have distinct identities (Buddhism challenged Brahmanism).
\end{itemize}
\end{QuickRevision}

\section{P. V. Kane: ``Ought'' vs ``Is''}

\textbf{P. V. Kane} (author of \emph{History of Dharmashastra}) says Indology is more about \textbf{what India ought to be} than \textbf{what India is}.

\begin{ExampleBox}
The \textbf{Grihya Sutra} (householder duties) describes what an \textbf{ideal woman} ought to be --- looking after the household, caring for the elderly and the husband. But what India is today: women are in the job market, equally achievement-oriented after feminism. So there is a large gap between the textual ``ought'' and the actual ``is''.
\end{ExampleBox}

\begin{QuickRevision}
\begin{itemize}
\item P. V. Kane: Indology describes ``ought to be'' more than ``is'' (Grihya Sutra's ideal woman vs feminism).
\end{itemize}
\end{QuickRevision}

\section{Acceptance of the Aryan Invasion Theory}

Indology \textbf{starts with the acceptance of the Aryan invasion theory}. In doing so it subscribes to \textbf{British Orientalism} --- to how the British wanted India to be seen, as part of a power equation designed to justify colonial supremacy.

\begin{QuickRevision}
\begin{itemize}
\item Indology starts by accepting the Aryan invasion theory = subscribing to British Orientalism.
\end{itemize}
\end{QuickRevision}

\section{Edward Said \& Inden: Propaganda, Not a Discipline}

\textbf{Edward Said} (a post-colonial scholar) and \textbf{Inden} argue that Indology was \textbf{less a discipline and more a propaganda}. Post-colonial studies deal with what happened to colonies after liberation.

\begin{KeyConcept}
Indology gives a \textbf{fixed, timeless and spaceless view} of Indian society --- with no regional variation, let alone historical change. India is presented as static (Vedic, Hindu, Aryan, with no social change), which helped the British convince the indigenous people that they had come to bring change.
\end{KeyConcept}

\begin{CriticalPoint}
The nationalist Indologists' subscription to the Aryan theory is itself a subscription to the power equation designed by the British --- a discourse that gave elevated status to outsiders and a demeaning status to the indigenous population. \textbf{Ashis Nandy} (\emph{Colonization of Mind}) says even the way we criticise the British is designed by them.
\end{CriticalPoint}

\begin{QuickRevision}
\begin{itemize}
\item Edward Said \& Inden: Indology = propaganda, giving a fixed, timeless, spaceless view of India.
\item Ashis Nandy, \emph{Colonization of Mind}; the discourse is one of power (Foucault: knowledge is power).
\end{itemize}
\end{QuickRevision}

\section{Carol Upadhyay: Brahmanical \& Savarna Outlook}

\begin{CriticalPoint}
\textbf{Carol Upadhyay} (the same scholar who traced Ghurye's sociological imagination to three factors): the Indological approach described Hindu customs and rituals as Indian society, making the discipline not only \textbf{politically conservative} but also \textbf{Brahmanical and Savarna} in its outlook.
\end{CriticalPoint}

\begin{QuickRevision}
\begin{itemize}
\item Carol Upadhyay: Indology made the discipline politically conservative, Brahmanical and Savarna.
\end{itemize}
\end{QuickRevision}

\section{Lack of Objectivity \& Empiricism}

Indology \textbf{lacks objectivity and empiricism}. It is based on text and interpretation, and that interpretation is simply assumed to be Indian culture --- there is no objectivity in the whole research.

\begin{QuickRevision}
\begin{itemize}
\item Indology lacks objectivity and empiricism (text + interpretation assumed as Indian culture).
\end{itemize}
\end{QuickRevision}

\section{T. K. Oommen: A West European Notion of Nation}

\begin{CriticalPoint}
\textbf{T. K. Oommen} (Professor Emeritus, JNU) says Ghurye held a \textbf{notion of nation that is utterly West European}. When Ghurye speaks of the ``assimilation of non-Aryans into Vedic Aryans'', he is using colonial language --- diffusionism and evolutionism are themselves Western theories. The very idea of nation Ghurye held is a colonial construction.
\end{CriticalPoint}

\begin{QuickRevision}
\begin{itemize}
\item T. K. Oommen: Ghurye's notion of nation is West European; ``assimilation'' is colonial language.
\end{itemize}
\end{QuickRevision}

\section{The Brahmin Reflexivity Problem}

\textbf{D. P. Mukherjee} believed a true Indian sociologist must be well-versed in Sanskrit; for him only Ghurye was the ``Indian sociologist'', the rest mere ``sociologists of India''.

\begin{CriticalPoint}
But whether Ghurye, D. P. Mukherjee, M. N. Srinivas or Andre Beteille --- \textbf{all four syllabus scholars are Brahmin men}. As Brahmins they subscribe to theories in which Brahmins are superior; there is \textbf{no reflexivity} in their sociology. Mukherjee's position neglects the perspective of those who did not belong to the Vedic Aryan culture.
\end{CriticalPoint}

\begin{QuickRevision}
\begin{itemize}
\item All four syllabus scholars are Brahmin men --- no reflexivity about their social location.
\item D. P. Mukherjee: only Ghurye is the ``Indian sociologist'' (Sanskrit-based).
\end{itemize}
\end{QuickRevision}

\section{Conclusion: Dumont \& Pocock}

In conclusion, Indology established the discourse of \textbf{cultural studies} in India, out of which Indian sociology emerged. \textbf{Dumont and Pocock} believe that Indian sociology lies at the \textbf{intersection of Indology and sociology}.

\begin{QuickRevision}
\begin{itemize}
\item Indology established cultural studies; Indian sociology emerged from it.
\item Dumont \& Pocock: Indian sociology lies at the intersection of Indology and sociology.
\end{itemize}
\end{QuickRevision}

\section{The Contribution of Indology to Indian Sociology (PYQ)}

\begin{PYQLink}
``How did Ghurye's Indological Perspective help in the contribution to Indian Sociology?'' --- the question asks for the \textbf{contribution}, not merely an explanation of Indology. First explain the perspective, then Ghurye as Indologist, and finally the contribution in 3--4 lines.
\end{PYQLink}

The positive contributions of the Indological perspective to Indian sociology:
\begin{itemize}
\item \textbf{It started the discourse of cultural studies} in India --- people began writing on Indian/Hindu culture, Hinduization and Aryanization.
\item \textbf{It started the comparative study of culture} --- between Vedic Aryans and European Aryans, between caste and race (Varna = colour), through philology/Indo-European. Comparative study is the essence of sociology --- we compare and generalise.
\item \textbf{It established the subject matter of Indian sociology = the study of social history} --- Ghurye laid the foundation stone (the evolution of caste, comparison with its past).
\item \textbf{It studied Hindu culture and how it established solidarity/integration all over India} --- the same collective consciousness from Vedic Hindu Aryan culture; all became Sanskritized/Hinduized/integrated. This is why some call Ghurye the ``Indian Durkheim''.
\item \textbf{It linked Western and Indian sociology} --- Dumont and Pocock: Indian sociology lies at the confluence of Indology and Western sociology.
\end{itemize}

\begin{CriticalPoint}
Later concepts --- Sanskritization, little/great tradition, purity/pollution --- all trace back to Ghurye's debate on Hinduization. (Note: Ghurye on caste is different --- there he says caste developed from race and occupation stabilised it, but only in Hindustan proper (North/Central India); in South India caste and race were not synonymous.)
\end{CriticalPoint}

\begin{DataFact}
Study reference: NCERT Class 11, Chapter 5 ``Indian Sociologists'', pages 85--90 (for Ghurye).
\end{DataFact}

\begin{QuickRevision}
\begin{itemize}
\item Indology's contributions: cultural studies, comparative study of culture, social history as subject matter, Hindu culture/solidarity, link to Western sociology.
\item Dumont \& Pocock: Indian sociology = confluence of Indology and Western sociology.
\end{itemize}
\end{QuickRevision}
